\section*{РОЗДІЛ 6. РИЗИКИ, ОБМЕЖЕННЯ ТА ЕТИЧНІ АСПЕКТИ}

Упровадження систем на основі великих мовних моделей у практики наукового адміністрування супроводжується низкою ризиків і обмежень, які мають як технічну, так і соціально-правову природу. Для проєкту ResearchIQ ці аспекти є визначальними, оскільки помилки в інтерпретації бібліометричних даних, некоректні рекомендації чи порушення режимів обробки персональних даних можуть призводити до управлінських рішень із довготривалими наслідками.

\subsection*{6.1. Технічні ризики та обмеження}

\subsubsection*{6.1.1. Проблема галюцинацій}

Одним із ключових технічних ризиків при використанні LLM є явище \emph{галюцинацій} (hallucinations), тобто генерації неправдивої інформації, яка подається у формі переконливого та стилістично коректного тексту. Для задач наукового адміністрування це особливо критично, адже навіть поодинокі хибні твердження щодо публікацій, цитувань чи авторства здатні спотворити аналітику та підірвати довіру користувачів до системи.

Емпіричні оцінки свідчать, що проблема має суттєвий масштаб: у порівняльних дослідженнях точності референсів, згенерованих LLM у задачах наукового письма, повідомляються значні частки некоректних або вигаданих посилань \cite{Chelli2024HallucinationRates}. У складніших постановках із багатокроковим міркуванням (multi-hop reasoning) ризик фактологічних помилок і накопичення похибки зростає, що робить критичними процедури верифікації та контроль provenance.

У бібліометричному контексті галюцинації проявляються у кількох типових формах (табл.~\ref{tab:hallucinations-biblio}): від створення неіснуючих публікацій і вигаданих метрик до помилкової атрибуції цитат і спотворення метаданих.

\begin{longtable}{@{}p{0.22\textwidth}p{0.46\textwidth}p{0.22\textwidth}@{}}
\caption{Типові прояви галюцинацій у бібліометричних задачах}\label{tab:hallucinations-biblio}\\
\toprule
\textbf{Тип} & \textbf{Приклад} & \textbf{Ризик} \\
\midrule
\textbf{Вигадані публікації} & Посилання на неіснуючу статтю з правдоподібними авторами/журналом & Високий \\
\textbf{Неправильна атрибуція} & Приписування цитати іншому автору & Середній \\
\textbf{Спотворені метадані} & Неправильний рік, журнал, кількість цитувань & Середній \\
\textbf{Вигадані метрики} & Неіснуючий H-індекс, citation count & Високий \\
\textbf{Вигадані тренди} & «Тренди», сформовані на основі галюцинованих даних & Високий \\
\bottomrule
\end{longtable}

Відповідно, архітектура ResearchIQ має передбачати системні механізми зниження цього ризику. До практичних стратегій мітигації належать: (i) верифікація фактів і метаданих за авторитетними зовнішніми джерелами (зокрема Crossref/OpenAlex), (ii) використання оцінювання впевненості (confidence scoring) та явне маркування невизначених фрагментів, (iii) надання посилань на першоджерела для відтворюваності перевірки, а також (iv) залучення користувача до остаточного затвердження критично важливих висновків (human-in-the-loop).

\subsubsection*{6.1.2. Залежність від зовнішніх API}

Окремою групою технічних ризиків є залежність системи від зовнішніх сервісів і програмних інтерфейсів (API), які використовуються для ідентифікації авторів та отримання бібліографічних даних (наприклад, ORCID, OpenAlex). Ризики тут пов'язані з доступністю (downtime), обмеженнями на швидкість запитів (rate limits) та еволюцією самих API, зокрема їхнім застаріванням або зміною схем (API deprecation), що потребує постійного моніторингу та адаптації інтеграцій.

Для зменшення залежності доцільно застосовувати локальне кешування з керованим часом життя записів (TTL, Time-To-Live), механізми \emph{graceful degradation} (часткова працездатність системи при відмові одного з джерел), а також підтримку кількох альтернативних постачальників даних.

\subsubsection*{6.1.3. Масштабованість і продуктивність}

Питання масштабованості для ResearchIQ визначається як обмеженнями контексту LLM, так і ресурсною вартістю обробки запитів. Зокрема, великі профілі (понад 1000 публікацій) можуть перевищувати контекстні ліміти моделей, а режим потокової видачі (streaming) для інтерактивних відповідей вимагає оптимізації затримок. Додатково, при використанні комерційних API істотною стає вартість одного запиту (cost per query), що безпосередньо впливає на економіку експлуатації.

Практичними підходами до розв'язання цих обмежень є розбиття даних на блоки (chunking) та поетапне узагальнення (summarization), багаторівнева обробка (tiered approach), а також гібридна архітектура, де частина аналітики виконується детермінованими алгоритмами, а LLM застосовується переважно для інтерпретації та генерації пояснень.

\subsubsection*{6.1.4. Операційні політики: rate limits, кешування, вартість}

Для виробничого використання критичними є політики \textbf{rate limiting}, повторів (retries) та ідемпотентності (idempotency) при зверненні до зовнішніх джерел. Практично це означає: (i) керований backoff при помилках/лімітах, (ii) кешування результатів запитів із диференційованими TTL для різних джерел, (iii) журналювання кожного tool call для подальшого аудиту та відтворення.

Окремим аспектом є \textbf{економіка запиту} (cost per query) і контроль витрат (LLM/API/storage). У межах цієї роботи такі оцінки розглядаються як необхідний компонент наступного етапу оцінювання (\textit{Evidence gap}: відсутній повний cost breakdown для типових сценаріїв і масштабів).

\subsection*{6.2. Алгоритмічна упередженість}

\subsubsection*{6.2.1. Джерела упередженості}

Алгоритмічна упередженість є системним ризиком для будь-яких моделей, що навчаються на великих масивах текстових даних. Оскільки навчальні корпуси відображають наявні соціальні та інституційні асиметрії, LLM можуть як відтворювати, так і підсилювати упередження, формуючи нерівномірні рекомендації або інтерпретації. У контексті наукової аналітики це проявляється, зокрема, як англоцентричність, географічна упередженість і \emph{disciplinary bias}, що здатні впливати на ранжування робіт, тематичні «карти» та порівняння дослідників.

\subsubsection*{6.2.2. 3-D Dependable AI Framework}

Для системного управління цими ризиками доцільно застосовувати рамкові підходи на кшталт \emph{3-D Dependable AI Framework} (Design / Development / Deployment), який розглядає упередженість як проблему, що має бути контрольованою на всіх етапах життєвого циклу: від проєктування вимог і даних до тестування, моніторингу та операційної експлуатації.

\subsubsection*{6.2.3. Мітигація в ResearchIQ}

У межах ResearchIQ першочерговими є механізми прозорості (пояснення джерел і логіки рекомендацій), можливість користувача коригувати параметри та результати (user override), а також підтримка багатомовності для зменшення ефектів англоцентричності. Подальший розвиток системи може включати формалізований аудит упередженості (bias audit), балансування рекомендацій і звітність за метриками справедливості (fairness evaluation reports) для різних груп користувачів та дисциплін.

\subsection*{6.3. Регуляторні та правові аспекти}

Регуляторне середовище для систем ШІ динамічно змінюється, а вимоги до обробки персональних даних і прозорості алгоритмів стають дедалі жорсткішими. Для ResearchIQ ключовим є узгодження практик обробки даних із принципами законності, мінімізації та обмеження мети, а також забезпечення підзвітності.

\subsubsection*{6.3.1. GDPR та захист персональних даних}

У межах права ЄС орієнтиром слугують позиції Європейської ради із захисту даних: зокрема, \textbf{EDPB Opinion 28/2024} \cite{EDPB2024Opinion28} формулює вимоги та застереження щодо застосування ШІ-систем у контексті GDPR, акцентуючи на прозорості, правових підставах обробки й управлінні ризиками.

\subsubsection*{6.3.2. EU AI Act}

Додатковий контур вимог задає \textbf{Regulation (EU) 2024/1689} \cite{EU2024AIAct}, відома як EU AI Act, яка визначає підхід до класифікації ризиків ШІ-систем та встановлює обов'язки для постачальників і користувачів залежно від класу ризику.

\subsubsection*{6.3.3. Українське законодавство}

Станом на 2025 рік Україна не має спеціального закону про ШІ, однак застосовними залишаються норми щодо захисту персональних даних та загальні вимоги інформаційної безпеки, а також євроінтеграційні зобов'язання, які поступово наближають регуляторні підходи до практик ЄС.

\subsection*{6.4. Академічна доброчесність}

Використання інструментів генеративного ШІ в академічному середовищі піднімає питання доброчесності та належного розкриття внеску автоматизованих систем. Хоча ResearchIQ позиціонується як інструмент підтримки адміністративних та аналітичних процесів, а не генерації наукового контенту, прозорість щодо способів отримання даних і формування висновків є необхідною для збереження довіри.

\subsubsection*{6.4.1. Політики університетів}

Політики університетів щодо використання GenAI істотно різняться: від повної заборони до дозволеного застосування за умови обов'язкового розкриття (disclosure) або встановлення правил «під завдання» (per assignment). Така варіативність означає потребу в гнучких налаштуваннях системи та вбудованих механізмах документування того, де саме інструменти ШІ були використані.

\subsubsection*{6.4.2. Політики видавців}

Провідні наукові видавці також формалізують правила використання інструментів ШІ, зокрема щодо авторства, відповідальності за зміст та прозорості методів [25, 26]. Для ResearchIQ це створює додаткову вимогу забезпечувати відтворюваність аналітики та можливість аудиту джерел.

\subsubsection*{6.4.3. Імплікації для ResearchIQ}

З практичної точки зору ResearchIQ має підтримувати принципи transparency та disclosure: надавати користувачам зрозумілу інформацію про використані джерела даних, межі застосовності отриманих результатів, а також про роль LLM у формуванні текстових пояснень і рекомендацій.

\subsection*{6.5. Принцип Human-in-the-Loop}

\subsubsection*{6.5.1. Концепція}

Підхід \emph{Human-in-the-Loop} (HITL) передбачає, що людина залишається в контурі прийняття рішень і може підтверджувати, коригувати або відхиляти результати роботи алгоритмів. У задачах наукового адміністрування це є важливим елементом відповідального застосування ШІ, оскільки дозволяє поєднати швидкість автоматизованої обробки з експертною оцінкою контексту.

\subsubsection*{6.5.2. Реалізація в ResearchIQ}

У ResearchIQ принцип HITL реалізується через інтерфейси, що підтримують перевірку критичних тверджень і редагування результатів перед їхнім використанням у звітності. Зокрема, користувач має мати можливість переглядати першоджерела, підтверджувати відповідність метаданих, а також задавати правила, за яких система зобов'язана запитувати підтвердження (наприклад, при низьких значеннях confidence score або при виявленні суперечностей між джерелами).

\subsubsection*{6.5.3. Підзвітність та журналювання (Audit Trail)}

Невід'ємним доповненням до HITL є забезпечення підзвітності через \emph{audit trail}: збереження історії запитів, версій джерел даних, параметрів обробки та внесених користувачем правок. Таке журналювання підвищує відтворюваність результатів, полегшує розбір інцидентів і створює основу для регулярних аудитів якості даних, безпеки та відповідності регуляторним вимогам.